% --- Metadados do trabalho ---
\newcommand{\universidade}{Universidade Federal do Pará (UFPA)}
\newcommand{\instituto}{Instituto de Ciências Exatas e Naturais (ICEN)}
\newcommand{\faculdade}{Faculdade de Ciência da Computação}
\newcommand{\titulo}{Resenha: O Algoritmo de Babai para o Isomorfismo de Grafos}
\newcommand{\disciplina}{Grafos}
\newcommand{\professor}{Dr. Roberto Samarone dos Santos Araujo}

% Lista de autores separada por vírgula
\newcommand{\autores}{ Alessandro Reali Lopes Silva }

% E-mails separados por vírgula, dentro de chaves para formatação uniforme
\newcommand{\emails}{alessandro.silva@icen.ufpa.br}


% --- Idioma, codificacao e fonte base ---
\usepackage[utf8]{inputenc}
\usepackage[T1]{fontenc}
\usepackage[brazilian]{babel}
\usepackage{mathptmx}

% --- Matematica e simbolos ---
\usepackage{amsmath, amssymb, amsfonts}
\usepackage{siunitx}
\sisetup{group-separator = {.}, output-decimal-marker = {,}}
% --- Midia, graficos e codigo ---
\usepackage{graphicx}
\usepackage{xcolor}
\usepackage{listings}
\usepackage[font=small, labelfont=bf]{caption}
\usepackage{tikz}
\usepackage{pgfplots}
\pgfplotsset{compat=1.18}
\usepackage{booktabs}
\renewcommand{\arraystretch}{1.5}

% --- Layout e espacamento ---
\usepackage[a4paper, top=3cm, bottom=2cm, left=3cm, right=2cm]{geometry}
\usepackage[nopatch=footnote]{microtype}

\usepackage{titlesec}
\titlespacing*{\section}{0pt}{10pt}{10pt}
\titlespacing*{\subsection}{0pt}{8pt}{4pt}

\usepackage{parskip}

% --- Links e navegacao ---
\usepackage[colorlinks=true, linkcolor=black, urlcolor=blue]{hyperref}

% --- Estilo do Código (Listings) ---
\definecolor{codecomment}{rgb}{0.3, 0.3, 0.3} % Cinza escuro para comentários
\definecolor{codekeyword}{rgb}{0.0, 0.0, 0.0} % Preto negrito para keywords
\definecolor{codestring}{rgb}{0.2, 0.2, 0.2}  % Quase preto para strings

\lstset{
    basicstyle=\ttfamily\small,           % Fonte mono (Courier/Consolas)
    commentstyle=\itshape\color{codecomment},
    keywordstyle=\bfseries\color{codekeyword},
    stringstyle=\color{codestring},
    numbers=left,                         % Números à esquerda para referência
    numberstyle=\tiny\color{gray},
    stepnumber=1,
    numbersep=8pt,
    backgroundcolor=\color{white},
    showspaces=false,
    showstringspaces=false,
    showtabs=false,
    frame=lines,                          % Linha em cima e embaixo (estilo SBC)
    rulecolor=\color{black},
    tabsize=2,
    captionpos=b,                         % Legenda embaixo
    breaklines=true,                      % Quebra de linha automática
    breakatwhitespace=false,
    escapeinside={\%*}{*)},               % Para usar LaTeX dentro do código
    inputencoding=utf8,
    extendedchars=true,
    literate={á}{{\'a}}1 {é}{{\'e}}1 {í}{{\'i}}1 {ó}{{\'o}}1 {ú}{{\'u}}1 {ã}{{\~a}}1 {õ}{{\~o}}1 {ç}{{\c{c}}}1
}

% --- Macros e personalizacoes gerais ---
\newcommand{\divisoria}{\noindent\rule{\textwidth}{0.4pt}}
\newcommand{\sinal}[1]{\mathcal{S}(#1)}
\newcommand{\fourier}{\mathcal{F}}
\newcommand{\hz}{\text{Hz}}
\newcommand{\khz}{\text{kHz}}
\newcommand{\mhz}{\text{MHz}}
\newcommand{\ghz}{\text{GHz}}
\newcommand{\db}{\text{dB}}
\newcommand{\s}{\text{s}}
\newcommand{\ms}{\text{ms}}
\newcommand{\us}{\text{ }\mu\text{s}}
\newcommand{\bps}{\text{bps}}
\newcommand{\kbps}{\text{Kbps}}
\newcommand{\mbps}{\text{Mbps}}